%% =============================================================
%%  chapters/09_conclusion.tex
%% =============================================================
\chapter{Conclusion and Future Work}
\label{ch:conclusion}

\section{Summary of Contributions}
\label{sec:summary-contributions}
This thesis presented a multi-task deep learning architecture that fuses
camera images and seat-embedded pressure heatmaps to jointly estimate
occupancy, passenger class, weight, and pose for automotive cabin
monitoring. Key contributions include a synthetic pressure heatmap
generation pipeline based on human body simulation, a Sim2Real validation
procedure quantifying the domain gap between real and synthetic heatmaps,
a camera-pressure synchronisation pipeline, and an empirical comparison of
fusion strategies through ablation studies.
%% TODO: Refine this summary once final results are available, and add
%% any additional contributions (e.g. released dataset or code).

\section{Answers to Research Questions}
\label{sec:answers-research-questions}
Revisiting the research questions posed in
Section~\ref{sec:research-questions}: regarding RQ1, the ablation study in
Section~\ref{sec:ablation-study} indicates \todo{summarise whether and by
how much fusion improved each task}; regarding RQ2, the Sim2Real
validation in Section~\ref{sec:sim2real-validation} shows \todo{summarise
the measured statistical and domain-gap findings}; regarding RQ3, the
comparison of fusion alternatives in
Section~\ref{sec:fusion-alternatives-rejected} demonstrates \todo{state
which fusion strategy was found preferable and why}; regarding RQ4, the
Sim2Real transfer evaluation in
Section~\ref{sec:sim2real-transfer-evaluation} shows \todo{summarise the
effect of mixed real+synthetic training on real-only test performance}.
%% TODO: Replace all placeholders above with concrete, results-grounded
%% answers once Chapter 8 is finalised.

\section{Limitations}
\label{sec:limitations}
The real pressure and camera dataset collected for this thesis is limited
in subject diversity and recording duration due to practical laboratory
constraints, and the synthetic data generation pipeline relies on
simplifying assumptions in the human body and seat contact simulation
that may not capture all real-world seating behaviours. Additionally, the
evaluation was conducted offline on recorded data rather than on embedded
target hardware, so real-time performance under production constraints
remains to be confirmed.
%% TODO: Add any further limitations identified during experimentation,
%% e.g. specific failure modes discussed in Section~\ref{sec:discussion-failure-cases}.

\section{Future Work}
\label{sec:future-work}
Future work could extend the real data collection campaign to a larger
and more diverse subject population, incorporate additional sensing
modalities such as depth or seatbelt tension sensors, and evaluate the
proposed architecture on embedded automotive hardware to validate
real-time feasibility under the non-functional requirements defined in
Section~\ref{sec:non-functional-requirements}. Further investigation into
domain adaptation techniques could also help further close the Sim2Real
gap identified in Section~\ref{sec:domain-gap-assessment}.
%% TODO: Prioritise future work items based on supervisor feedback and
%% remaining time budget.

\section{Closing Remarks}
\label{sec:closing-remarks}
The fusion of camera and pressure sensor data investigated in this
thesis represents a step toward more robust, privacy-conscious cabin
monitoring systems capable of supporting the next generation of adaptive
vehicle safety features. The methodology and findings presented here are
intended to inform both future academic research and practical system
design in this area.
%% TODO: Personalise closing remarks once the thesis defence date and
%% final findings are confirmed.
